% !TeX root = ../main.tex
\section{Einleitung}
\label{sec:einleitung}

\subsection{Ausgangslage und Ziel}
Schwimmer messen sich an nationalen Meisterschaften, an internationalen
Wettkampfserien sowie an Welt- und Europameisterschaften. Ranglisten zeigen,
wer wie schnell schwimmt. Sie zeigen aber nicht, wer wem an Wettkämpfen
begegnet und wie sich daraus ein Netz des Schwimmsports bildet. Genau diesen
Blick ermöglicht die soziale Netzwerkanalyse. Wir verbinden Athleten, die am
selben Veranstaltungstag eine Bestleistung erzielten, und untersuchen die
Struktur dieses Netzes.

Grundlage sind die TOP100-Listen von swimrankings.net mit den persönlichen
Bestleistungen von über 13\,000 Athleten aus 182 Nationen. Erhoben wurden sie
mit einem Scraper aus einem früheren Projekt (Kapitel~\ref{sec:daten}). Ziel
ist zu verstehen, wie national oder international der Schwimmsport vernetzt
ist, welche Athleten und Veranstaltungen das Netz prägen und ob sich künftige
Begegnungen vorhersagen lassen.

\subsection{Fragestellungen}
Daraus leiten wir fünf Fragen ab.
\begin{enumerate}
  \item Wie stark ist das Netz national geprägt, und welche Gruppen bilden
  sich (Kapitel~\ref{sec:reduktion} und~\ref{sec:communities})?
  \item Welche Athleten und Nationalteams sind zentral, und sind stärkere
  Schwimmer zentraler (Kapitel~\ref{sec:zentralitaet})?
  \item Was hält das Netz zusammen, einzelne Athleten oder grosse
  Veranstaltungen (Kapitel~\ref{sec:zusammenhalt})?
  \item Mit wem verbinden sich Athleten bevorzugt
  (Kapitel~\ref{sec:assortativitaet})?
  \item Lassen sich künftige Begegnungen aus der Netzwerkstruktur vorhersagen
  (Kapitel~\ref{sec:link_prediction})?
\end{enumerate}

Die Untersuchung beschränkt sich auf die gespeicherten Bestleistungen der
TOP100-Listen mit Stand 2025. Eine Kante bedeutet, dass zwei Athleten am selben Veranstaltungstag eine ihrer Karrierebestzeiten erzielten. Sie belegt weder eine
persönliche Bekanntschaft noch ein direktes Rennen gegeneinander. Die Netze
fassen alle Jahre zusammen, nur die Link Prediction trennt zwei Zeiträume.
Die verwendeten Verfahren und Masse folgen, wo nicht anders angegeben, dem
Skript von \textcite{henninger2024}.
Kapitel~\ref{sec:daten} und~\ref{sec:modellierung} beschreiben Daten und
Modellierung, die Kapitel~\ref{sec:reduktion} bis~\ref{sec:link_prediction}
die Analysen. Kapitel~\ref{sec:schluss} beantwortet die Fragen und gibt einen
Ausblick.
